\scanpage{90}

\uline{Application} \quad Soit $X$ schéma projectif régulier sur
un corps $k$ \inserted{sép. clos}, $Y$ une section hyperplane
régulière, $F$ un faisceau de $l$-\deleted{\unsure{$l$|tor}}torsion loc.
constant sur $X$ \deleted{[$l$ premier à la car. de $k$]}.

On suppose \uline{$\dim X=n$}. Considérons $U=X-Y$ et

\diagram{\resizebox{\linewidth}{!}{\begin{tikzcd}[ampersand replacement=\&,column sep=small,row sep=1ex]
H^{i-1}(U,F)\arrow[r] \& H_Y^{\unsure{i|i-1}}(X,F)\arrow[r] \& H^i(X,F)\arrow[r] \& H^i(U,F)\arrow[r] \& H_Y^{i+1}(X,F)\\
\& \rotatebox{90}{$\simeq$} \& \& \& \rotatebox{90}{$\simeq$}\\
\& H^{i-2}(Y,F_Y\otimes\check T_{\deleted{\unsure{Y|l|\ell}}}) \& \& \& H^{\unsure{i+1|i-1}}(Y,F_Y\otimes\check T_{\deleted{\unsure{N|l|\ell}}})
\end{tikzcd}}}

et utilisons le fait que

\[ H^i(U,F)=0 \quad\text{si } i>n \]

On trouve le

\uline{Th de Lefschetz}

% vertical emphasis rule at left
\[
\begin{aligned}
\smash{\rule[-3ex]{.4pt}{5.8ex}}\qquad&\deleted{\unsure{H^i_Y(X,F)}}\\
&H^{i-2}(Y,F_Y\otimes\check T_l)\longrightarrow H^i(X,F)
\end{aligned}
\]

est \uline{surjectif} si $i\geqn+1$, \uline{bijectif} si $i\geq n+2$

\uline{Corollaire} \quad Supposons $X$ \uline{connexe} \inserted{\uline{régulier}} de dim $n$,
\llap{\smash{\rule[-1.8ex]{.4pt}{4.2ex}}\hspace{.7em}}alors $H^{2n}(X,\mu_N^{\otimes n})\simeq\mathbb Z/N\mathbb Z$ pour tt
entier $N$ premier à la car.

% archive sheet number: 89
